% GENERATED FILE. Edit canonical lessons, then run npm run generate:books.
% canonical-source-hash: fnv1a64:36d419ec15378b30
% canonical-lessons: MW-C195-baink, MW-C195-khato, MW-C195-karaj, MW-C195-baki, MW-C195-bachno

\chapter{At the Bank}
\label{ch:mw-at-the-bank}
\hlchaptermodality{195}

\begin{chapteropening}
\textbf{By the end of this chapter:} \emph{I can say what I keep and what I owe at the bank.}

Complete the last lesson of chapter 195: I can say what I keep and what I owe at the bank.
\end{chapteropening}

\section[\texorpdfstring{\mw{बैंक} (baiṅk)}{बैंक (baiṅk)}]{\mw{बैंक} (baiṅk) — a bank}

\label{lesson:MW-C195-baink}



\begin{quote}
Before the new one: say the Marwadi for a profit, then the Marwadi for credit.
\end{quote}

\subsection*{You'll want to know: \mw{बैंक}}
\textbf{\mw{बैंक}} — \emph{baiṅk} — \textquotedblleft{}a bank\textquotedblright{}.

\begin{quote}
\mw{बैंक} \mw{बाजार} \mw{कनै} \mw{है।}
\end{quote}

\begin{quote}
\emph{baiṅk bājār kanai hai.}
\end{quote}

\begin{quote}
The bank is near the market.
\end{quote}

A masculine noun.

\begin{cousinweb}
\textbf{\mw{बैंक}} (\emph{baiṅk}) is borrowed from English \emph{bank}.
\end{cousinweb}

\begin{grammarlens}[title={What you've built}]
One more word for reading notices and messages.
\end{grammarlens}

\subsection*{Your turn}
\begin{itemize}
\raggedright
  \item \emph{Say it:} \emph{baiṅk}
  \item \emph{Say it:} \emph{baiṅk}, once more
  \item \emph{Say it:} \emph{baiṅk bājār kanai hai.} — The bank is near the market.
  \item \emph{Recall:} say the Marwadi for a profit, then the Marwadi for credit, then say \emph{baiṅk} again
\end{itemize}

\begin{tcolorbox}[breakable,colback=teal!4,colframe=teal!35!black,title={Before you move on}]
What does \mw{बैंक} mean? (\textquotedblleft{}A bank\textquotedblright{}.) Say it once more.
\end{tcolorbox}
\section[\texorpdfstring{\mw{खातो} (khāto)}{खातो (khāto)}]{\mw{खातो} (khāto) — a bank account}

\label{lesson:MW-C195-khato}



\begin{quote}
Before the new one: say the Marwadi for credit, then the Marwadi for a bank.
\end{quote}

\subsection*{You'll want to know: \mw{खातो}}
\textbf{\mw{खातो}} — \emph{khāto} — \textquotedblleft{}a bank account\textquotedblright{}.

\begin{quote}
\mw{बैंक} \mw{मांय} \mw{म्हारो} \mw{खातो} \mw{है।}
\end{quote}

\begin{quote}
\emph{baiṅk mā̃y mhāro khāto hai.}
\end{quote}

\begin{quote}
I have an account at the bank.
\end{quote}

A masculine noun.

\begin{cousinweb}
\textbf{\mw{खातो}} (\emph{khāto}) is filed by CDIAL under Sanskrit \emph{kṣatra}, \textquotedblleft{}might, rule\textquotedblright{}, with Hindi \emph{khātā} and Sindhi \emph{khāto}, \textquotedblleft{}account book\textquotedblright{} (CDIAL 3648). CDIAL gives no Marwari form in this entry.
\end{cousinweb}

\begin{grammarlens}[title={What you've built}]
One more word for reading notices and messages.
\end{grammarlens}

\subsection*{Your turn}
\begin{itemize}
\raggedright
  \item \emph{Say it:} \emph{khāto}
  \item \emph{Say it:} \emph{khāto}, once more
  \item \emph{Say it:} \emph{baiṅk mā̃y mhāro khāto hai.} — I have an account at the bank.
  \item \emph{Recall:} say the Marwadi for credit, then the Marwadi for a bank, then say \emph{khāto} again
\end{itemize}

\begin{tcolorbox}[breakable,colback=teal!4,colframe=teal!35!black,title={Before you move on}]
What does \mw{खातो} mean? (\textquotedblleft{}A bank account\textquotedblright{}.) Say it once more.
\end{tcolorbox}
\section[\texorpdfstring{\mw{करज} (karaj)}{करज (karaj)}]{\mw{करज} (karaj) — a debt}

\label{lesson:MW-C195-karaj}



\begin{quote}
Before the new one: say the Marwadi for a bank, then the Marwadi for a bank account.
\end{quote}

\subsection*{You'll want to know: \mw{करज}}
\textbf{\mw{करज}} — \emph{karaj} — \textquotedblleft{}a debt\textquotedblright{}.

\begin{quote}
\mw{बाप} \mw{माथै} \mw{घणो} \mw{करज} \mw{है।}
\end{quote}

\begin{quote}
\emph{bāp māthai ghaṇo karaj hai.}
\end{quote}

\begin{quote}
Father has a large debt.
\end{quote}

A masculine noun.

\begin{cousinweb}
\textbf{\mw{करज}} (\emph{karaj}) is usually traced, through Persian, to Arabic \emph{qarḍ}, \textquotedblleft{}a loan\textquotedblright{}.
\end{cousinweb}

\begin{grammarlens}[title={What you've built}]
One more word for reading notices and messages.
\end{grammarlens}

\subsection*{Your turn}
\begin{itemize}
\raggedright
  \item \emph{Say it:} \emph{karaj}
  \item \emph{Say it:} \emph{karaj}, once more
  \item \emph{Say it:} \emph{bāp māthai ghaṇo karaj hai.} — Father has a large debt.
  \item \emph{Recall:} say the Marwadi for a bank, then the Marwadi for a bank account, then say \emph{karaj} again
\end{itemize}

\begin{tcolorbox}[breakable,colback=teal!4,colframe=teal!35!black,title={Before you move on}]
What does \mw{करज} mean? (\textquotedblleft{}A debt\textquotedblright{}.) Say it once more.
\end{tcolorbox}
\section[\texorpdfstring{\mw{बाकी} (bākī)}{बाकी (bākī)}]{\mw{बाकी} (bākī) — remaining; still owed}

\label{lesson:MW-C195-baki}



\begin{quote}
Before the new one: say the Marwadi for a bank account, then the Marwadi for a debt.
\end{quote}

\subsection*{You'll want to know: \mw{बाकी}}
\textbf{\mw{बाकी}} — \emph{bākī} — \textquotedblleft{}remaining; still owed\textquotedblright{}.

\begin{quote}
\mw{दस} \mw{रुपिया} \mw{बाकी} \mw{है।}
\end{quote}

\begin{quote}
\emph{das rupiyā bākī hai.}
\end{quote}

\begin{quote}
Ten rupees are still owed.
\end{quote}

\begin{cousinweb}
\textbf{\mw{बाकी}} (\emph{bākī}) is usually traced, through Persian, to Arabic \emph{bāqī}, \textquotedblleft{}remaining\textquotedblright{}.
\end{cousinweb}

\begin{grammarlens}[title={What you've built}]
One more word for reading notices and messages.
\end{grammarlens}

\subsection*{Your turn}
\begin{itemize}
\raggedright
  \item \emph{Say it:} \emph{bākī}
  \item \emph{Say it:} \emph{bākī}, once more
  \item \emph{Say it:} \emph{das rupiyā bākī hai.} — Ten rupees are still owed.
  \item \emph{Recall:} say the Marwadi for a bank account, then the Marwadi for a debt, then say \emph{bākī} again
\end{itemize}

\begin{tcolorbox}[breakable,colback=teal!4,colframe=teal!35!black,title={Before you move on}]
What does \mw{बाकी} mean? (\textquotedblleft{}Remaining; still owed\textquotedblright{}.) Say it once more.
\end{tcolorbox}
\section[\texorpdfstring{\mw{बचणो} (bachṇo)}{बचणो (bachṇo)}]{\mw{बचणो} (bachṇo) — to be left over; to escape}

\label{lesson:MW-C195-bachno}



\begin{quote}
Before the new one: say the Marwadi for a debt, then the Marwadi for remaining.
\end{quote}

\subsection*{You'll want to know: \mw{बचणो}}
\textbf{\mw{बचणो}} — \emph{bachṇo} — \textquotedblleft{}to be left over; to escape\textquotedblright{}.

\begin{quote}
\mw{म्हारी} \mw{तनखा} \mw{सूं} \mw{सो} \mw{रुपिया} \mw{बचै} \mw{है।}
\end{quote}

\begin{quote}
\emph{mhārī tankhā sū̃ so rupiyā bachai hai.}
\end{quote}

\begin{quote}
A hundred rupees are left over from my salary.
\end{quote}

Past (it …): \mw{वो} \mw{सिक्को} \mw{बच्यो}, \emph{vo sikko bachyo}.

\begin{cousinweb}
\textbf{\mw{बचणो}} (\emph{bachṇo}) is inherited from Sanskrit \emph{vañcati}, \textquotedblleft{}moves; escapes from\textquotedblright{}; CDIAL lists Old Marwari \emph{vacaï}, \textquotedblleft{}is saved\textquotedblright{} (CDIAL 11208). \textbf{\mw{बचावणो}}, \textquotedblleft{}to save\textquotedblright{}, is its partner.
\end{cousinweb}

\begin{grammarlens}[title={What you've built}]
One more word for reading notices and messages.
\end{grammarlens}

\subsection*{Your turn}
\begin{itemize}
\raggedright
  \item \emph{Say it:} \emph{bachṇo}
  \item \emph{Say it:} \emph{bachṇo}, once more
  \item \emph{Say it:} \emph{mhārī tankhā sū̃ so rupiyā bachai hai.} — A hundred rupees are left over from my salary.
  \item \emph{Recall:} say the Marwadi for a debt, then the Marwadi for remaining, then say \emph{bachṇo} again
\end{itemize}

\begin{tcolorbox}[breakable,colback=teal!4,colframe=teal!35!black,title={Before you move on}]
What does \mw{बचणो} mean? (\textquotedblleft{}To be left over; to escape\textquotedblright{}.) Say it once more.
\end{tcolorbox}
